\documentclass[../DoAn.tex]{subfiles}
\begin{document}

Chương 4 đã trình bày cách hiện thực các module của hệ thống. Chương này trình bày kết quả kiểm thử và đánh giá hệ thống ở hai khía cạnh. Phần đầu (mục 5.1--5.2) đánh giá phần \textit{sản phẩm}: 28 trường hợp kiểm thử chức năng và đánh giá hybrid parser so với full-LLM parser. Phần sau (mục 5.3) trình bày \textit{đóng góp nghiên cứu chính} của khóa luận: thí nghiệm so sánh định lượng ba phương pháp matching CV--JD trên 400 cặp dữ liệu công khai, với kiểm định thống kê đầy đủ để trả lời hai câu hỏi nghiên cứu RQ1 và RQ2.

\section{Kiểm thử chức năng}

Hệ thống được kiểm thử theo phương pháp kiểm thử hộp đen, tập trung vào việc xác minh các chức năng hoạt động đúng theo yêu cầu. Ngoài ra, các unit test bằng Jest và Testing Library được viết cho các component React quan trọng.

\subsection{Kiểm thử module xác thực}

Bảng \ref{tab:test_auth} trình bày kết quả kiểm thử cho module xác thực.

\begin{table}[H]
\centering
\caption{Kết quả kiểm thử module xác thực}
\label{tab:test_auth}
\resizebox{\textwidth}{!}{%
\begin{tabular}{|c|p{5cm}|p{4cm}|c|}
\hline
\textbf{TC} & \textbf{Mô tả} & \textbf{Kết quả mong đợi} & \textbf{Trạng thái} \\ \hline
TC01 & Đăng ký với email hợp lệ & Tạo tài khoản thành công & Đạt \\ \hline
TC02 & Đăng ký với email đã tồn tại & Thông báo lỗi email trùng & Đạt \\ \hline
TC03 & Đăng ký với mật khẩu yếu & Thông báo yêu cầu mật khẩu mạnh & Đạt \\ \hline
TC04 & Đăng nhập bằng email/mật khẩu đúng & Đăng nhập thành công, chuyển hướng & Đạt \\ \hline
TC05 & Đăng nhập bằng username (admin alias) & Đăng nhập thành công & Đạt \\ \hline
TC06 & Đăng nhập với mật khẩu sai & Thông báo lỗi xác thực & Đạt \\ \hline
TC07 & Đăng nhập bằng Google OAuth & Chuyển hướng Google, callback thành công & Đạt \\ \hline
TC08 & Đăng nhập bằng LinkedIn OAuth & Chuyển hướng LinkedIn, callback thành công & Đạt \\ \hline
TC09 & Truy cập trang bảo vệ khi chưa đăng nhập & Chuyển hướng đến trang đăng nhập & Đạt \\ \hline
TC10 & Truy cập editor khi là khách & Cho phép truy cập (phiên khách) & Đạt \\ \hline
\end{tabular}%
}
\end{table}

\subsection{Kiểm thử module tải lên và phân tích CV}

Bảng \ref{tab:test_upload} trình bày kết quả kiểm thử cho module tải lên và phân tích CV.

\begin{table}[H]
\centering
\caption{Kết quả kiểm thử module tải lên và phân tích CV}
\label{tab:test_upload}
\resizebox{\textwidth}{!}{%
\begin{tabular}{|c|p{5cm}|p{4cm}|c|}
\hline
\textbf{TC} & \textbf{Mô tả} & \textbf{Kết quả mong đợi} & \textbf{Trạng thái} \\ \hline
TC12 & Upload file PDF hợp lệ ($<$10MB) & Trích xuất văn bản và phân tích thành công & Đạt \\ \hline
TC13 & Upload file DOCX hợp lệ & Trích xuất và phân tích thành công & Đạt \\ \hline
TC14 & Upload file không hỗ trợ (.jpg) & Thông báo lỗi định dạng & Đạt \\ \hline
TC15 & Upload file quá 10MB & Thông báo vượt kích thước cho phép & Đạt \\ \hline
TC16 & Upload CV tiếng Việt & Nhận diện ngôn ngữ VI, phân tích đúng & Đạt \\ \hline
TC17 & Upload CV tiếng Anh & Nhận diện ngôn ngữ EN, phân tích đúng & Đạt \\ \hline
TC18 & Upload file không phải CV & \texttt{possibility\_score} thấp ($<$30) & Đạt \\ \hline
TC19 & Phân tích khi OpenAI API lỗi & Chuyển sang fallback parser & Đạt \\ \hline
\end{tabular}%
}
\end{table}

\subsection{Kiểm thử trình soạn thảo CV}

Bảng \ref{tab:test_editor} trình bày kết quả kiểm thử trình soạn thảo CV.

\begin{table}[H]
\centering
\caption{Kết quả kiểm thử trình soạn thảo CV}
\label{tab:test_editor}
\resizebox{\textwidth}{!}{%
\begin{tabular}{|c|p{5cm}|p{4cm}|c|}
\hline
\textbf{TC} & \textbf{Mô tả} & \textbf{Kết quả mong đợi} & \textbf{Trạng thái} \\ \hline
TC20 & Chỉnh sửa thông tin liên hệ & Cập nhật và hiển thị trên preview & Đạt \\ \hline
TC21 & Thêm mục kinh nghiệm mới & Thêm thành công, hiển thị trên preview & Đạt \\ \hline
TC22 & Xoá mục kinh nghiệm & Xoá thành công với xác nhận & Đạt \\ \hline
TC23 & Kéo thả sắp xếp thứ tự & Thứ tự cập nhật trên cả editor và preview & Đạt \\ \hline
TC24 & Thêm/xoá kỹ năng dạng tag & Tag thêm/xoá đúng & Đạt \\ \hline
TC25 & Auto-save sau chỉnh sửa & Dữ liệu lưu tự động sau 2 giây & Đạt \\ \hline
TC26 & Preview cập nhật real-time & Preview phản ánh ngay khi sửa & Đạt \\ \hline
TC27 & Thêm section tuỳ chỉnh & Section mới hiển thị đúng & Đạt \\ \hline
\end{tabular}%
}
\end{table}

\subsection{Kiểm thử module xuất CV}

Bảng \ref{tab:test_export} trình bày kết quả kiểm thử cho module xuất CV.

\begin{table}[H]
\centering
\caption{Kết quả kiểm thử module xuất CV}
\label{tab:test_export}
\begin{tabular}{|c|p{5cm}|p{4cm}|c|}
\hline
\textbf{TC} & \textbf{Mô tả} & \textbf{Kết quả mong đợi} & \textbf{Trạng thái} \\ \hline
TC28 & Xuất PDF & Mở print dialog, layout đúng & Đạt \\ \hline
TC29 & Xuất DOCX & Tải file .docx, mở được bằng Word & Đạt \\ \hline
\end{tabular}
\end{table}

Tổng cộng 28 trường hợp kiểm thử đều đạt. Chi tiết kịch bản và đầu vào của từng test case được liệt kê trong Phụ lục \ref{tab:test_auth_apx}.

\section{Đánh giá Hybrid Parser so với Full-LLM Parser}
\label{sec:hybrid_eval}

Mục này so sánh phương pháp \textit{hybrid parser} (regex + LLM cho kinh nghiệm) với \textit{full-LLM parser} (dùng LLM cho toàn bộ CV) trên 20 CV mẫu (10 tiếng Việt, 10 tiếng Anh) với cấu trúc đa dạng. Bảng \ref{tab:parser_comparison} trình bày kết quả so sánh.

\begin{table}[H]
\centering
\caption{So sánh hiệu năng Hybrid Parser và Full-LLM Parser}
\label{tab:parser_comparison}
\begin{tabular}{|l|c|c|}
\hline
\textbf{Tiêu chí} & \textbf{Full-LLM} & \textbf{Hybrid} \\ \hline
Thời gian trung bình (giây) & 8,5 & 4,2 \\ \hline
Token sử dụng trung bình & 3\,200 & 1\,400 \\ \hline
Chi phí ước tính / CV (USD) & 0,0048 & 0,0021 \\ \hline
Độ chính xác thông tin liên hệ & 95\% & 93\% \\ \hline
Độ chính xác kinh nghiệm & 90\% & 88\% \\ \hline
Độ chính xác kỹ năng & 88\% & 85\% \\ \hline
Tỷ lệ giảm token & -- & 56\% \\ \hline
\end{tabular}
\end{table}

Phương pháp hybrid giảm khoảng 56\% lượng token API sử dụng và giảm hơn một nửa thời gian xử lý so với full-LLM. Mức giảm nhẹ về độ chính xác (2--3\%) ở các trường thông tin liên hệ và kỹ năng là chấp nhận được, do các trường này có cấu trúc tương đối cố định và regex xử lý đủ tốt. Đối với kinh nghiệm làm việc, mức chênh lệch 2\% (90\% vs 88\%) chủ yếu đến từ các CV có cấu trúc không chuẩn -- nơi LLM có lợi thế hơn trong việc hiểu ngữ cảnh.

\section{Thí nghiệm so sánh ba phương pháp matching CV--JD}
\label{sec:experiment}

Đây là đóng góp nghiên cứu chính của khóa luận, trả lời hai câu hỏi nghiên cứu RQ1 (cách nào cho xếp hạng tốt nhất) và RQ2 (cách nào đáng đồng tiền nhất khi đưa vào sử dụng thực tế).

\subsection{Thiết kế thí nghiệm}

\textbf{Dữ liệu}. Bộ dữ liệu gồm 400 cặp (CV, mô tả công việc) lấy từ Kaggle Resume Dataset và Kaggle LinkedIn Job Postings -- hai bộ dữ liệu công khai. Ground truth được xây dựng theo \textit{Strategy A}: dùng category-mapping (24 nhóm nghề) để gắn nhãn nhị phân:

\begin{itemize}
\item \textit{Positive} (label = 1): CV và JD thuộc cùng nhóm nghề. 200 cặp.
\item \textit{Negative} (label = 0): CV và JD thuộc khác nhóm nghề. 200 cặp.
\end{itemize}

Bộ dữ liệu được xáo trộn bằng \texttt{seed = 42} để tái lập, sau đó cả ba phương pháp được chạy lần lượt trên cùng 400 cặp này.

\textbf{Phương pháp}. Ba phương pháp matching được mô tả ở §\ref{sec:ir} (Chương 2): TF-IDF (cách 1), Embedding với \texttt{text-embedding-3-small} (cách 2), và LLM với \texttt{gpt-4o-mini} kèm \texttt{temperature = 0}, \texttt{seed = 42} (cách 3).

\textbf{Chỉ số đánh giá}. Để trả lời RQ1, khóa luận dùng bốn chỉ số xếp hạng chuẩn (xem §\ref{sec:metrics}): Precision@10, nDCG@10, MRR, ROC-AUC. Để trả lời RQ2, khóa luận đo độ trễ (p50, p95) và chi phí thực (USD trên 1\,000 cặp).

\textbf{Kiểm định thống kê}. Cho mỗi cặp phương pháp, sai số mỗi phương pháp được tính là $|\text{label} - \text{score}/100|^2$. Sau đó dùng \textit{paired Wilcoxon signed-rank test} để kiểm tra xem khác biệt sai số giữa hai phương pháp có ý nghĩa thống kê ($p < 0,05$) hay không. Khoảng tin cậy 95\% của nDCG@10 được tính bằng \textit{bootstrap} với 10\,000 lần lấy mẫu và \texttt{seed = 42}.

\subsection{Kết quả định lượng}

Bảng \ref{tab:rq1_main} trình bày kết quả các chỉ số xếp hạng cho ba phương pháp.

\begin{table}[H]
\centering
\caption{Kết quả các chỉ số xếp hạng trên 400 cặp (RQ1)}
\label{tab:rq1_main}
\begin{tabular}{|l|c|c|c|}
\hline
\textbf{Chỉ số} & \textbf{TF-IDF} & \textbf{Embedding} & \textbf{LLM} \\ \hline
Precision@5  & 1,000 & 1,000 & 1,000 \\ \hline
Precision@10 & 1,000 & 1,000 & 1,000 \\ \hline
nDCG@10      & 1,000 & 1,000 & 1,000 \\ \hline
MRR          & 1,000 & 1,000 & 1,000 \\ \hline
ROC-AUC      & 0,755 & \textbf{0,766} & 0,755 \\ \hline
F1 (ngưỡng median) & 0,695 & 0,695 & \textbf{0,717} \\ \hline
\end{tabular}
\end{table}

Các chỉ số xếp hạng top-$k$ (Precision@5, P@10, nDCG@10, MRR) đều đạt 1,0 ở cả ba phương pháp -- do bộ dữ liệu có sự phân tách giữa hai lớp tương đối rõ ràng, dẫn đến cả ba phương pháp đều xếp đúng các cặp dương ở top. Chỉ số phân biệt được rõ nhất là \textbf{ROC-AUC}: Embedding đạt 0,766 -- cao nhất trong ba phương pháp, vượt cả TF-IDF (0,755) và LLM (0,755).

Bảng \ref{tab:rq2_cost} trình bày độ trễ và chi phí cho 1\,000 cặp.

\begin{table}[H]
\centering
\caption{Độ trễ và chi phí cho 1\,000 cặp (RQ2)}
\label{tab:rq2_cost}
\begin{tabular}{|l|c|c|c|}
\hline
\textbf{Chỉ số} & \textbf{TF-IDF} & \textbf{Embedding} & \textbf{LLM} \\ \hline
Độ trễ p50 (ms) & 1   & 0\,$^{*}$ & 2\,527 \\ \hline
Độ trễ p95 (ms) & 3   & 1\,$^{*}$ & 5\,214 \\ \hline
Chi phí (USD / 1\,000 cặp) & 0 & 0\,$^{*}$ & 0,28 \\ \hline
\end{tabular}
\\\smallskip
\footnotesize{$^{*}$ Embedding chạy với cache đã pre-compute. Chi phí pre-compute toàn bộ corpus 676 văn bản chỉ tốn khoảng 0,024 USD và chạy một lần.}
\end{table}

LLM tốn 0,28 USD cho 1\,000 cặp và chậm hơn Embedding khoảng \textit{5\,000 lần} (2\,527 ms vs 0 ms khi đã cache). TF-IDF nhanh và miễn phí (1 ms, \$0), tuy nhiên chất lượng (ROC-AUC) thấp hơn Embedding.

\subsection{Kiểm định thống kê}

Bảng \ref{tab:wilcoxon} trình bày kết quả kiểm định paired Wilcoxon signed-rank giữa các cặp phương pháp, dựa trên sai số bình phương từng cặp.

\begin{table}[H]
\centering
\caption{Kiểm định paired Wilcoxon signed-rank trên sai số bình phương (n = 400)}
\label{tab:wilcoxon}
\begin{tabular}{|l|c|c|c|}
\hline
\textbf{So sánh} & \textbf{Statistic} & \textbf{p-value} & \textbf{Phương pháp sai số thấp hơn} \\ \hline
Embedding vs LLM    & 33\,402 & 0,0038       & \textbf{Embedding} \\ \hline
Embedding vs TF-IDF & 20\,551 & $2{,}9 \times 10^{-17}$ & \textbf{Embedding} \\ \hline
LLM vs TF-IDF       & 24\,135 & $5{,}2 \times 10^{-12}$ & LLM \\ \hline
\end{tabular}
\end{table}

Tất cả ba cặp so sánh đều có $p$-value $< 0{,}05$, nghĩa là khác biệt sai số là \textbf{có ý nghĩa thống kê}. Cụ thể:

\begin{itemize}
\item Embedding có sai số trung bình \textbf{thấp hơn LLM} một cách đáng kể ($p = 0{,}0038$).
\item Embedding có sai số thấp hơn nhiều so với TF-IDF ($p < 10^{-16}$).
\item LLM có sai số thấp hơn TF-IDF ($p < 10^{-11}$).
\end{itemize}

Thứ tự sai số tăng dần: \textbf{Embedding $<$ LLM $<$ TF-IDF}. Khoảng tin cậy 95\% bootstrap cho ROC-AUC của Embedding là [0,72; 0,81] -- không bao hàm giá trị 0,5 (ngẫu nhiên), khẳng định kết quả không phải do may rủi.

\subsection{Trả lời các câu hỏi nghiên cứu}

\textbf{RQ1 -- Cách nào cho xếp hạng tốt nhất?} Trên dữ liệu 400 cặp Kaggle, cách \textit{so sánh nghĩa} (Embedding với \texttt{text-embedding-3-small}) cho kết quả phân biệt tốt nhất với ROC-AUC = 0,766, vượt cả TF-IDF (0,755) và LLM (0,755). Kiểm định Wilcoxon xác nhận khác biệt giữa Embedding và hai phương pháp còn lại đều có ý nghĩa thống kê ($p < 0{,}05$). Cả ba phương pháp đều đạt $P@10 = 1{,}0$ và $\text{nDCG}@10 = 1{,}0$ vì bộ dữ liệu có hai lớp phân tách rõ ở top-$k$.

\textbf{RQ2 -- Cách nào đáng đồng tiền nhất?} Kết quả ở Bảng \ref{tab:rq2_cost} cho thấy \textit{Embedding} là cách hiệu quả chi phí nhất: đạt chất lượng cao nhất với chi phí gần như bằng 0 sau khi cache. LLM cho chất lượng tương đương TF-IDF (ROC-AUC = 0,755) nhưng chậm 5\,000 lần và tốn \$0,28/1\,000 cặp. TF-IDF có ưu thế \textit{khả năng triển khai dễ dàng} (không cần API key, không phụ thuộc dịch vụ ngoài), phù hợp khi hệ thống không thể gọi API ngoài hoặc cần xử lý hàng triệu cặp.

Tóm lại, với bài toán so khớp CV--JD ở quy mô vài trăm đến vài chục nghìn cặp, \textbf{Embedding là lựa chọn cân bằng tốt nhất giữa chất lượng và chi phí}.

\section{So sánh với các giải pháp hiện có}

Bảng \ref{tab:full_comparison} so sánh chi tiết hệ thống của khóa luận với bốn giải pháp phổ biến trên thị trường.

\begin{table}[H]
\centering
\caption{So sánh chi tiết hệ thống với các giải pháp hiện có}
\label{tab:full_comparison}
\resizebox{\textwidth}{!}{%
\begin{tabular}{|p{3.5cm}|c|c|c|c|c|}
\hline
\textbf{Tiêu chí} & \textbf{Canva} & \textbf{Zety} & \textbf{Sovren} & \textbf{LinkedIn} & \textbf{Khóa luận} \\ \hline
Upload CV có sẵn & Không & Không & Có & Không & Có \\ \hline
Phân tích bằng AI/LLM & Không & Hạn chế & NLP truyền thống & Không & LLM (GPT) \\ \hline
Soạn thảo tương tác & Template & Form & Không & Profile & Section + DnD \\ \hline
Xem trước real-time & Có & Có & Không & Không & Có \\ \hline
Xuất PDF & Có & Có & Không & Hạn chế & Có \\ \hline
Xuất DOCX & Không & Có & Không & Không & Có \\ \hline
Hỗ trợ tiếng Việt & Hạn chế & Không & Hạn chế & Hạn chế & Đầy đủ \\ \hline
Đa phương thức xác thực & Email & Email & API key & LinkedIn & Email + Google + LinkedIn \\ \hline
Phiên khách & Không & Không & Không & Không & Có \\ \hline
Auto-save & Có & Có & -- & Có & Có \\ \hline
Mã nguồn mở & Không & Không & Không & Không & Có \\ \hline
Chi phí & Freemium & Trả phí & Trả phí cao & Miễn phí & Miễn phí \\ \hline
\end{tabular}%
}
\end{table}

Hệ thống của khóa luận có lợi thế rõ ở các điểm: (i) kết hợp được cả tải lên CV, phân tích bằng AI, và soạn thảo tương tác trong cùng một hệ thống; (ii) hỗ trợ tiếng Việt đầy đủ; (iii) cho phép phiên khách dùng mà không cần đăng ký; (iv) mã nguồn mở cho phép cộng đồng đóng góp và tùy chỉnh.

\section{Hạn chế và bài học rút ra}

Mặc dù kết quả thí nghiệm khá rõ ràng, khóa luận có một số hạn chế cần lưu ý:

\textbf{Hạn chế về ground truth.} Strategy A dán nhãn theo category-mapping là \textit{silver-standard}: hai cặp cùng nhóm nghề không nhất thiết là phù hợp 100\%, và hai cặp khác nhóm có thể vẫn liên quan (ví dụ frontend developer và full-stack developer). Việc dùng ground truth do người chấm điểm tay (gold-standard) sẽ cho kết quả tin cậy hơn nhưng tốn nhiều công sức.

\textbf{Hạn chế về phạm vi dữ liệu.} Bộ dữ liệu chỉ gồm 400 cặp tiếng Anh, chưa bao gồm CV tiếng Việt. Kết quả có thể khác trên dữ liệu tiếng Việt do mô hình embedding chủ yếu được huấn luyện trên tiếng Anh.

\textbf{Hạn chế về quy mô.} Số cặp (n = 400) đủ để có kiểm định thống kê có ý nghĩa, nhưng vẫn nhỏ so với các benchmark IR chuẩn (thường vài chục nghìn cặp).

\textbf{Bài học.} Thí nghiệm cho thấy LLM không phải lúc nào cũng là phương pháp tốt nhất: trên bài toán so khớp CV--JD, một mô hình embedding nhỏ hơn rất nhiều (\texttt{text-embedding-3-small}) cho chất lượng cao hơn và chi phí thấp hơn nhiều lần. Điều này nhắc nhở rằng việc \textit{chọn đúng công cụ} cho từng bài toán quan trọng hơn việc dùng công cụ phức tạp nhất có sẵn.

\end{document}
